% !TeX program = pdflatex
% Compilar desde la raíz del proyecto TFM.
\documentclass[11pt]{article}
\usepackage[utf8]{inputenc}
\usepackage[T1]{fontenc}
\usepackage[spanish,es-noshorthands]{babel}
\usepackage{lmodern,amsmath,amssymb,bm,tikz,hyperref}
\usepackage[a3paper,landscape,margin=18mm]{geometry}
\usetikzlibrary{arrows.meta,shapes.geometric,babel}
\definecolor{azulTFM}{HTML}{155C8A}
\definecolor{verdeTFM}{HTML}{2D7F5E}
\definecolor{naranjaTFM}{HTML}{D97706}
\newcommand{\vect}[1]{\bm{#1}}
\pagestyle{empty}
\begin{document}
\centering
{\LARGE\bfseries Resolución numérica de la llama\par}
\vspace{3mm}
{\large Esquema global de KFLAME: del estado inicial a la aceptación\par}
\vspace{5mm}
\input{.local/research/thesis/figuras/flujograma_resolucion_numerica_llama.tikz.tex}

\end{document}
